\section{Reference Frames and Sensor Models}
\label{sec:reference_frames}

% ----------------------------------------------------------------
% TARGET: ~1.5 pages. Two frames + two sensor models.
% ----------------------------------------------------------------

\subsection{The World Frame and the Body Frame}
\label{sec:frames:definition}

We define two right-handed Cartesian coordinate systems
(Fig.~\ref{fig:frames}):

\begin{itemize}
  \item \textbf{World frame} $\mathcal{W} = \{X_w, Y_w, Z_w\}$. A fixed,
        Earth-referenced frame. We adopt the East-North-Up (ENU) convention:
        $X_w$ points East, $Y_w$ points North, and $Z_w$ points upward
        (opposite to gravity).
  \item \textbf{Body frame} $\mathcal{B} = \{X_b, Y_b, Z_b\}$. A frame
        rigidly attached to the sensor. Its axes rotate with the device.
\end{itemize}

At any instant, the orientation of $\mathcal{B}$ relative to $\mathcal{W}$ can
be described by a rotation. Euler angles (roll~$\phi$, pitch~$\theta$,
yaw~$\psi$) are intuitive but suffer from \emph{gimbal lock}---a singularity
at $\theta = \pm 90°$ where one degree of freedom is lost. Quaternions
avoid this problem entirely and are the representation we adopt in this work
(Section~\ref{sec:quaternions}).

\begin{figure}[t]
  \centering
  \includegraphics[width=\columnwidth]{Yaw_Axis_Corrected.png}
  \caption{The three Euler rotation axes illustrated on a rigid body.
           \textbf{Roll}~($\phi$) is the rotation about the longitudinal axis,
           \textbf{Pitch}~($\theta$) about the lateral axis, and
           \textbf{Yaw}~($\psi$) about the vertical axis.
           Any orientation of the body frame~$\mathcal{B}$ relative to the
           world frame~$\mathcal{W}$ can be decomposed into a sequence of
           these three rotations, though the representation is subject to
           gimbal lock at $\theta = \pm 90°$.}
  \label{fig:frames}
\end{figure}

\subsection{Accelerometer: Angle from Gravity Projection}
\label{sec:frames:accel}

\subsubsection{Sensor model}

An ideal tri-axial accelerometer measures the \emph{specific force}
experienced by the sensor, i.e.\ the difference between the true linear
acceleration and gravity, expressed in the body frame:

\begin{equation}
  \mathbf{a}_b = R^{-1}\!\bigl(\mathbf{a}_\text{true} - \mathbf{g}\bigr)
               + \mathbf{b}_a + \mathbf{n}_a
  \label{eq:accel_model}
\end{equation}

\noindent where $R$ is the rotation matrix from $\mathcal{W}$ to $\mathcal{B}$,
$\mathbf{g} = [0,\, 0,\, -g]^\top$ is the gravity vector in $\mathcal{W}$
(with $g \approx 9.81~\text{m/s}^2$),
$\mathbf{b}_a$ is a slowly varying bias, and $\mathbf{n}_a$ is additive white
noise.

\subsubsection{Angle derivation under static conditions}

When the sensor is \textbf{stationary}
($\mathbf{a}_\text{true} = \mathbf{0}$), Eq.~\eqref{eq:accel_model}
simplifies to:

\begin{equation}
  \mathbf{a}_b \approx R^{-1}(-\mathbf{g})
  \label{eq:accel_static}
\end{equation}

\noindent That is, the accelerometer reads only the projection of the gravity
vector onto the three body axes. Because gravity is a known, constant vector,
we can \emph{invert} this projection to recover the tilt of the sensor. We now
derive the exact formulas using rotation matrices.

\subsubsection{Formal derivation via rotation matrices}
\label{sec:frames:accel:derivation}

We model the accelerometer as a three-dimensional coordinate system
$(x, y, z)$ that rotates with respect to a fixed Earth frame where gravity
always points upward: $\mathbf{g} = [0,\; 0,\; g]^\top$ (with $z$ pointing up
and $g \approx 9.81~\text{m/s}^2$).

\paragraph{Step 1: Rotation matrices.}
We define the elementary Tait-Bryan rotations. Roll ($\phi$) is a rotation
about the $x$-axis, and pitch ($\theta$) is a rotation about the $y$-axis:

\begin{equation}
  R_x(\phi) = \begin{bmatrix}
    1 & 0 & 0 \\
    0 & \cos\phi & \sin\phi \\
    0 & -\sin\phi & \cos\phi
  \end{bmatrix}
  \label{eq:Rx}
\end{equation}

\begin{equation}
  R_y(\theta) = \begin{bmatrix}
    \cos\theta & 0 & -\sin\theta \\
    0 & 1 & 0 \\
    \sin\theta & 0 & \cos\theta
  \end{bmatrix}
  \label{eq:Ry}
\end{equation}

When the sensor is at rest, the measured acceleration vector $\mathbf{a}_b =
(A_x, A_y, A_z)^\top$ is the gravity vector transformed by the sensor's
orientation:

\begin{equation}
  \begin{bmatrix} A_x \\ A_y \\ A_z \end{bmatrix}
  = R_y(\theta)\;R_x(\phi)\;\begin{bmatrix} 0 \\ 0 \\ g \end{bmatrix}
  \label{eq:accel_rotation}
\end{equation}

\paragraph{Step 2: Matrix multiplication.}
We first apply $R_x(\phi)$ to the gravity vector:

\begin{equation}
  R_x(\phi)\begin{bmatrix} 0 \\ 0 \\ g \end{bmatrix}
  = \begin{bmatrix} 0 \\ g\sin\phi \\ g\cos\phi \end{bmatrix}
  \label{eq:Rx_g}
\end{equation}

\noindent Then we apply $R_y(\theta)$ to this result:

\begin{equation}
  R_y(\theta)\begin{bmatrix} 0 \\ g\sin\phi \\ g\cos\phi \end{bmatrix}
  = \begin{bmatrix}
      -g\sin\theta\cos\phi \\[2pt]
      g\sin\phi \\[2pt]
      g\cos\theta\cos\phi
    \end{bmatrix}
  \label{eq:Ry_Rx_g}
\end{equation}

\noindent Therefore, the three measured components are:

\begin{align}
  A_x &= -g\,\sin\theta\,\cos\phi \label{eq:Ax} \\
  A_y &= \phantom{-}g\,\sin\phi \label{eq:Ay} \\
  A_z &= \phantom{-}g\,\cos\theta\,\cos\phi \label{eq:Az}
\end{align}

\paragraph{Step 3: Solving for roll ($\phi$).}
We divide Eq.~\eqref{eq:Ay} by Eq.~\eqref{eq:Az}:

\begin{equation}
  \frac{A_y}{A_z}
  = \frac{g\,\sin\phi}{g\,\cos\theta\,\cos\phi}
  \label{eq:roll_ratio}
\end{equation}

\noindent However, to obtain a result that is independent of $\theta$, we use
the more robust form. Computing $\sqrt{A_x^2 + A_z^2}$:

\begin{equation}
  \sqrt{A_x^2 + A_z^2}
  = g\,\cos\phi\,\sqrt{\sin^2\!\theta + \cos^2\!\theta}
  = g\,\cos\phi
  \label{eq:Axz_norm}
\end{equation}

\noindent Dividing $A_y$ by this result:

\begin{equation}
  \frac{A_y}{\sqrt{A_x^2 + A_z^2}}
  = \frac{g\,\sin\phi}{g\,\cos\phi}
  = \tan\phi
\end{equation}

\noindent Therefore:

\begin{equation}
  \boxed{\;\phi = \arctan\!\left(
    \frac{A_y}{\sqrt{A_x^2 + A_z^2}}
  \right)\;}
  \label{eq:roll_accel}
\end{equation}

\paragraph{Step 4: Solving for pitch ($\theta$).}
From Eq.~\eqref{eq:Ax}, we have $A_x = -g\,\sin\theta\,\cos\phi$. We need to
relate $A_x$ to the remaining components. Squaring and summing $A_y$ and
$A_z$:

\begin{align}
  A_y^2 + A_z^2
    &= g^2\sin^2\!\phi + g^2\cos^2\!\theta\,\cos^2\!\phi \notag \\
    &= g^2\cos^2\!\phi\bigl(\tan^2\!\phi + \cos^2\!\theta\bigr)
  \label{eq:Ayz_sq}
\end{align}

\noindent For the standard Tait-Bryan decomposition where $|\theta| < 90°$,
the simpler path uses the identity $\sin^2\!\phi + \cos^2\!\phi = 1$.
Squaring Eqs.~\eqref{eq:Ay} and \eqref{eq:Az} directly:

\begin{equation}
  A_y^2 + A_z^2
  = g^2\bigl(\sin^2\!\phi + \cos^2\!\theta\,\cos^2\!\phi\bigr)
\end{equation}

\noindent Noting that the full norm satisfies
$A_x^2 + A_y^2 + A_z^2 = g^2$, we can write
$\sqrt{A_y^2 + A_z^2} = g\,\cos\theta\;$ (valid when $\cos\phi \approx 1$ or
by using the complete norm). Dividing $-A_x$ by this:

\begin{equation}
  \frac{-A_x}{\sqrt{A_y^2 + A_z^2}}
  = \frac{g\,\sin\theta}{g\,\cos\theta}
  = \tan\theta
\end{equation}

\noindent Therefore:

\begin{equation}
  \boxed{\;\theta = \arctan\!\left(
    \frac{-A_x}{\sqrt{A_y^2 + A_z^2}}
  \right)\;}
  \label{eq:pitch_accel}
\end{equation}

\paragraph{Summary of the derivation.}
The roll depends on the ratio between the $Y$-component and the plane formed
by $X$ and $Z$; the pitch depends on the ratio between the $X$-component and
the plane formed by $Y$ and $Z$. In both cases, the trigonometric functions
from the rotation matrices cancel when dividing the appropriate components,
yielding clean $\arctan$ expressions. Note that the \emph{yaw} ($\psi$)
cannot be recovered: a rotation about the vertical axis leaves the gravity
projection unchanged. Resolving yaw requires an additional sensor such as a
magnetometer.

\subsubsection{Limitation: sensitivity to dynamic motion}
\label{sec:frames:accel:limitation}

Equations~\eqref{eq:roll_accel}--\eqref{eq:pitch_accel} are exact only when
$\mathbf{a}_\text{true} = \mathbf{0}$. As soon as the sensor undergoes
\emph{any} acceleration---starting to move, changing direction, or
vibrating---the accelerometer reading becomes the superposition of gravity and
the dynamic acceleration:

\begin{equation}
  \mathbf{a}_b = R^{-1}\!\bigl(\mathbf{a}_\text{true} - \mathbf{g}\bigr)
  \label{eq:accel_dynamic}
\end{equation}

\noindent In this regime, applying
Eqs.~\eqref{eq:roll_accel}--\eqref{eq:pitch_accel} yields angles that are
dominated by the dynamic term, rendering the estimate unreliable. The
accelerometer is, in effect, a \textbf{low-frequency sensor for attitude}: its
angle estimate is trustworthy when the sensor is at rest or moving at constant
velocity, but \emph{loses all precision during transients} (acceleration,
deceleration, vibration). Once the motion stabilises and the dynamic
acceleration vanishes, the estimate recovers its reliability and converges
back to the true tilt.

% ----- FIGURE: accelerometer behaviour -----
% TODO: Three-panel time-series plot:
%   Panel 1: sensor at rest       → angle accurate, slightly noisy
%   Panel 2: sensor accelerating  → angle wildly wrong
%   Panel 3: sensor at rest again → angle recovers
%
% \begin{figure}[t]
%   \centering
%   \includegraphics[width=\columnwidth]{accel_behaviour.pdf}
%   \caption{Accelerometer-derived tilt angle during three phases:
%            \textbf{(a)}~stationary---accurate but noisy;
%            \textbf{(b)}~dynamic motion---dominated by linear acceleration,
%            estimate becomes unreliable;
%            \textbf{(c)}~stationary again---recovers to the true value.}
%   \label{fig:accel_behaviour}
% \end{figure}


\subsection{Gyroscope: Angle from Angular Velocity Integration}
\label{sec:frames:gyro}

\subsubsection{Sensor model}

A tri-axial gyroscope measures the instantaneous angular velocity of
$\mathcal{B}$ relative to $\mathcal{W}$, expressed in body-frame coordinates:

\begin{equation}
  \boldsymbol{\omega}_b = \boldsymbol{\omega}_\text{true}
                        + \mathbf{b}_\omega + \mathbf{n}_\omega
  \label{eq:gyro_model}
\end{equation}

\noindent where $\mathbf{b}_\omega$ is the gyroscope bias and
$\mathbf{n}_\omega$ is white noise.

\subsubsection{Angle derivation by integration}

Because the gyroscope measures angular \emph{velocity}, the angle traversed
around any axis is recovered by integration. For a single axis (e.g.\ roll),
the continuous-time relation is:

\begin{equation}
  \phi(t) = \phi(t_0) + \int_{t_0}^{t} \omega_x(\tau)\,d\tau
  \label{eq:gyro_angle_continuous}
\end{equation}

In practice, the sensor is sampled at a fixed rate $f_s = 1/\Delta t$, so the
integral is approximated by a discrete sum. Using the simplest (rectangular)
rule:

\begin{equation}
  \phi[n] = \phi[n-1] + \omega_x[n]\,\Delta t
  \label{eq:gyro_angle_discrete}
\end{equation}

\noindent and analogously for pitch and yaw. The key strength of this approach
is that the gyroscope is \textbf{insensitive to linear acceleration}: it
responds only to rotations, making it highly accurate during dynamic motion
when the accelerometer fails.

\subsubsection{Limitation: cumulative drift}
\label{sec:frames:gyro:drift}

Equation~\eqref{eq:gyro_angle_discrete} reveals the fundamental weakness of
pure gyroscope integration. At every time step, the measured angular velocity
carries an error $\epsilon[n] = b_\omega + n_\omega[n]$ composed of the
constant bias and a random noise sample. The integrated angle becomes:

\begin{equation}
  \phi[N] = \phi[0]
           + \underbrace{\sum_{n=1}^{N} \omega_\text{true}[n]\,\Delta t}%
             _{\text{true angle}}
           + \underbrace{N\,b_\omega\,\Delta t}%
             _{\text{bias drift}}
           + \underbrace{\sum_{n=1}^{N} n_\omega[n]\,\Delta t}%
             _{\text{random walk}}
  \label{eq:gyro_drift}
\end{equation}

Two error terms are visible:

\begin{enumerate}
  \item \textbf{Bias drift.} The constant bias $b_\omega$ accumulates linearly
        with each sample. After $N$ steps the error is $N\,b_\omega\,\Delta t
        = b_\omega\,T$, where $T$ is the total elapsed time. This term grows
        \emph{without bound}.
  \item \textbf{Random walk.} The noise samples $n_\omega[n]$ are uncorrelated
        with zero mean. Their cumulative effect has a standard deviation of
        $\sigma_\omega\sqrt{\Delta t \cdot T}$ (angular random walk), where
        $\sigma_\omega$ is the noise standard deviation. This grows as
        $\sqrt{T}$---much slower than the bias drift. A full derivation is
        given in Appendix~\ref{app:rw}.
\end{enumerate}

\subsubsection{Role of sampling frequency}

The magnitude of the discretisation error (the difference between the true
integral and the rectangular sum) depends on the time step $\Delta t$. For a
smooth signal $\omega(t)$, the local truncation error per step is
$\mathcal{O}(\Delta t^2)$, and the global accumulated error over a fixed
interval $T = N\Delta t$ is:

\begin{equation}
  \epsilon_\text{disc} = \mathcal{O}(N \cdot \Delta t^2)
                       = \mathcal{O}(\Delta t)
  \label{eq:disc_error}
\end{equation}

\noindent Therefore, \textbf{increasing the sampling frequency}
($\Delta t \to 0$) directly reduces the discretisation error. A sensor
sampling at 200~Hz accumulates roughly half the discretisation error of one
sampling at 100~Hz over the same time interval. However, no increase in
sampling rate can eliminate the bias drift, which depends on $b_\omega$ and the
total elapsed time, not on $\Delta t$.

\begin{figure}[t]
  \centering
  \includegraphics[width=\columnwidth]{gyro_drift_errors.pdf}
  \caption{Simulated gyroscope integration errors over 30~seconds
           ($b_\omega = 0.05$~deg/s, $\sigma_\omega = 0.1$~deg/s,
           $f_s = 100$~Hz).
           \textbf{(a)}~Bias drift grows linearly with time.
           \textbf{(b)}~Random walk is stochastic and grows as
           $\sigma_\omega\sqrt{\Delta t \cdot T}$ (dashed $1\sigma$ envelope,
           see Appendix~\ref{app:rw}).
           \textbf{(c)}~Combined error: the bias drift dominates at
           longer time scales.}
  \label{fig:gyro_drift}
\end{figure}

\begin{figure}[t]
  \centering
  \includegraphics[width=\columnwidth]{gyro_sampling_comparison.pdf}
  \caption{Effect of sampling frequency on gyroscope integration.
           \textbf{Top:} integrated angle at 50~Hz and 200~Hz compared to
           ground truth.
           \textbf{Bottom:} accumulated error. Higher sampling rates
           reduce discretisation error but cannot eliminate bias drift,
           which depends only on $b_\omega$ and elapsed time~$T$.}
  \label{fig:gyro_sampling}
\end{figure}

\subsection{Magnetometer: Heading from Earth's Magnetic Field}
\label{sec:frames:mag}

\subsubsection{Sensor model}

A tri-axial magnetometer measures the Earth's magnetic field vector expressed
in the body frame:

\begin{equation}
  \mathbf{m}_b = R^{-1}\,\mathbf{m}_w + \mathbf{b}_m + \mathbf{n}_m
  \label{eq:mag_model}
\end{equation}

\noindent where $\mathbf{m}_w$ is the local geomagnetic field vector in the
world frame, $R$ is the rotation matrix from $\mathcal{W}$ to $\mathcal{B}$,
$\mathbf{b}_m$ is a bias term caused by hard-iron distortions (nearby
ferromagnetic materials), and $\mathbf{n}_m$ is additive noise.

In the ENU convention, the geomagnetic field has a horizontal component
pointing approximately toward magnetic North (along $Y_w$) and a vertical
component whose sign and magnitude depend on the latitude (the
\emph{inclination} angle). At mid-latitudes, the field magnitude is typically
$25$--$65~\mu$T.

\subsubsection{Why the accelerometer cannot observe yaw}
\label{sec:frames:mag:why}

As derived in Section~\ref{sec:frames:accel:derivation}, the accelerometer
measures the projection of the gravity vector onto the body axes. Gravity is
a \emph{vertical} vector: $\mathbf{g} = [0,\;0,\;g]^\top$ in the world
frame. A rotation about the vertical axis (yaw, $\psi$) leaves this vector
\emph{invariant}:

\begin{equation}
  R_z(\psi)\,\mathbf{g}
  = \begin{bmatrix}
      \cos\psi & -\sin\psi & 0 \\
      \sin\psi & \cos\psi  & 0 \\
      0        & 0         & 1
    \end{bmatrix}
  \begin{bmatrix} 0 \\ 0 \\ g \end{bmatrix}
  = \begin{bmatrix} 0 \\ 0 \\ g \end{bmatrix}
  \label{eq:yaw_gravity_invariant}
\end{equation}

\noindent Regardless of the value of $\psi$, the accelerometer reading does
not change. Therefore, \textbf{yaw is unobservable from gravity alone}.

\paragraph{Important clarification: world yaw vs.\ body-axis rotation.}
The unobservable degree of freedom is always the rotation about the
\emph{world} vertical axis (the direction of~$\mathbf{g}$), not about the
body-frame $Z_b$~axis. When the device is held horizontally---as in a
drone---$Z_b$ is aligned with gravity, so a rotation about $Z_b$ coincides
with world yaw and is entirely invisible to the accelerometer. However, when
the device is tilted (e.g.\ a smartphone held upright), $Z_b$ is no longer
parallel to~$\mathbf{g}$. In that configuration, a rotation about $Z_b$
\emph{does} change the gravity projection and is therefore partially
observable. Nevertheless, there remains exactly \textbf{one unobservable
degree of freedom}---the rotation about the gravity vector---regardless of
device orientation. It is this world-frame yaw that the magnetometer must
resolve.

This is the fundamental reason why a two-sensor (accelerometer + gyroscope)
fusion pipeline cannot fully determine orientation: the gyroscope can track
yaw changes in the short term, but its estimate drifts without any absolute
reference to correct it.

\subsubsection{Heading derivation from the magnetometer}
\label{sec:frames:mag:heading}

The Earth's magnetic field, unlike gravity, has a significant \emph{horizontal}
component that changes direction with yaw rotations. To extract the heading
angle, the magnetometer reading must first be \emph{tilt-compensated} using
the known roll and pitch angles (obtained from the accelerometer or the
complementary filter). The procedure is:

\paragraph{Step 1: Tilt compensation.}
Rotate the raw magnetometer vector from the (tilted) body frame into the
horizontal plane using the inverse of the roll and pitch rotations:

\begin{equation}
  \begin{bmatrix} m_x' \\ m_y' \\ m_z' \end{bmatrix}
  = R_x(-\phi)\;R_y(-\theta)\;
  \begin{bmatrix} m_x \\ m_y \\ m_z \end{bmatrix}
  \label{eq:tilt_compensation}
\end{equation}

\noindent Expanding the matrix product, the tilt-compensated horizontal
components are:

\begin{align}
  m_x' &= m_x\cos\theta + m_y\sin\phi\sin\theta
         + m_z\cos\phi\sin\theta
  \label{eq:mx_comp} \\[4pt]
  m_y' &= m_y\cos\phi - m_z\sin\phi
  \label{eq:my_comp}
\end{align}

\paragraph{Step 2: Compute heading.}
The yaw angle is the bearing of the horizontal magnetic component:

\begin{equation}
  \boxed{\;\psi = \arctan2\!\bigl(-m_x',\; m_y'\bigr)\;}
  \label{eq:yaw_mag}
\end{equation}

\noindent The $\arctan2$ function returns a full-circle result in
$(-180°, 180°]$, resolving the quadrant ambiguity. Note that
$\psi = 0$ corresponds to magnetic North.

\subsubsection{Limitations of the magnetometer}
\label{sec:frames:mag:limitations}

\begin{itemize}
  \item \textbf{Soft- and hard-iron distortions.} Ferromagnetic materials
        near the sensor (e.g.\ a phone case with a magnetic clasp, a metal
        table) create constant offsets (hard-iron) or directional scaling
        (soft-iron) that corrupt the reading. Calibration---typically an
        ellipsoid fit---is required before use.
  \item \textbf{Low bandwidth.} The geomagnetic field changes slowly;
        consumer magnetometers are typically sampled at 50--100~Hz but the
        useful signal content is below $\sim$\!1~Hz. High-frequency heading
        changes are better tracked by the gyroscope.
  \item \textbf{Indoor interference.} Electronic devices, power cables, and
        reinforced-concrete structures generate local magnetic anomalies that
        can cause heading errors of tens of degrees.
\end{itemize}

\noindent Despite these limitations, the magnetometer is the \emph{only}
passive sensor available on a consumer device that provides an absolute
heading reference. Its role in the fusion pipeline is strictly analogous to
that of the accelerometer for tilt: it supplies a low-frequency, drift-free
yaw reference that is fused with the high-frequency gyroscope estimate via
the complementary filter.

\subsection{The Need for Fusion}
\label{sec:frames:fusion_need}

The preceding analysis reveals a natural complementarity among the three
sensors, summarised in Table~\ref{tab:sensor_comparison}:

\begin{itemize}
  \item The \textbf{accelerometer} provides an absolute tilt reference
        (gravity) that is \emph{bounded} in error---it never drifts---but is
        only valid when the sensor is at rest or in uniform motion.
        During transients it is unreliable; once the motion ends, it
        \emph{recovers}. It \emph{cannot} observe yaw.
  \item The \textbf{magnetometer} provides an absolute heading reference
        (Earth's magnetic field) that is also bounded in error and
        self-correcting, but is sensitive to local magnetic disturbances and
        has low bandwidth.
  \item The \textbf{gyroscope} tracks rotations on \emph{all three axes}
        faithfully during fast, dynamic motion, but its integrated angle
        \emph{drifts without bound} due to bias accumulation. It never
        self-corrects.
\end{itemize}

\noindent In the frequency domain, the accelerometer and magnetometer are
reliable \emph{low-frequency} sources (slow, steady orientation on
tilt and heading respectively), while the gyroscope is a reliable
\emph{high-frequency} source (fast changes on all axes). A well-designed
fusion algorithm should retain the low-frequency content of the absolute
sensors and the high-frequency content of the gyroscope, discarding the
noisy high-frequency accelerometer/magnetometer data and the drifting
low-frequency gyroscope data. This is precisely what the complementary
filter achieves (Section~\ref{sec:complementary_filter}).

\begin{table}[t]
  \centering
  \caption{Complementary characteristics of the three sensors for
           full attitude estimation.}
  \label{tab:sensor_comparison}
  \resizebox{\columnwidth}{!}{%
  \begin{tabular}{@{}lccc@{}}
    \toprule
    \textbf{Property} & \textbf{Accel.} & \textbf{Magnet.} & \textbf{Gyro.} \\
    \midrule
    Measures            & Specific force  & Magnetic field    & Angular vel.    \\
    Provides            & Roll, Pitch     & Yaw (heading)     & Roll, Pitch, Yaw \\
    Angle via           & Trigonometry    & $\arctan2$        & Integration     \\
    Low-freq.\ acc.     & High (gravity)  & High (mag.\ field)& Low (drift)     \\
    High-freq.\ acc.    & Low (dyn.\ acc.)& Low (noise)       & High            \\
    Absolute ref.       & Yes             & Yes               & No              \\
    Affected by motion  & Yes             & No                & No              \\
    Affected by env.    & No              & Yes (mag.\ anom.) & No              \\
    Drift               & None (bounded)  & None (bounded)    & Unbounded       \\
    Self-correcting     & Yes             & Yes               & No              \\
    \bottomrule
  \end{tabular}%
  }
\end{table}

% ----- FIGURE: sensor summary -----
% TODO: Side-by-side qualitative comparison:
%   Left panel  - "Accelerometer angle": noisy around true value,
%                  spikes during motion, returns after.
%   Right panel - "Gyroscope angle": smooth, tracks true value initially,
%                  but slowly diverges (drift).
%   Annotation:  arrow pointing to "fusion zone" where both contribute.
%
% \begin{figure}[t]
%   \centering
%   \includegraphics[width=\columnwidth]{sensor_summary.pdf}
%   \caption{Qualitative comparison of angle estimates.
%            \textbf{Left:} the accelerometer is accurate at rest but
%            corrupted during motion; it always recovers.
%            \textbf{Right:} the gyroscope is smooth and responsive but
%            drifts indefinitely. Fusing both yields the best of each.}
%   \label{fig:sensor_summary}
% \end{figure}
